% Generated prose for the post-result budget-paced follow-up (results/raw/followup-paced-v1).
% The table is produced by: python -m cc_cep.followup table --directory results/raw/followup-paced-v1
\subsection{Budget-paced follow-up (post-result)}\label{sec:followup}
After the main study we ran the paced policy of Table~\ref{tbl:policies} in 20 further runs: both fault types at both request rates, using the five v7 block seeds and the frozen v7 runtime. All 20 runs passed evidence verification and replay, giving 128 of 128 planned comparisons with v7 runs. The paced runs were collected about a week after v7 and paired by block ID only, so Table~\ref{tbl:followup} is a descriptive, time-separated comparison and is not pooled with the main study.

The paced policy spent 98--100\% of the usable allocation in every condition, compared with 29--30\% (aborts) and 93--94\% (delays) for the closed-loop policy. Abort exposure rose from 4.46 to 15.41 fraction-seconds at 32 requests/s and from 4.21 to 15.09 at 64 requests/s. Delay exposure rose by 0.98 and 0.76 fraction-seconds, and paced exposure exceeded open-loop exposure by 1.6--2.0 fraction-seconds in all four conditions. Paced exposure was within 0.62 fraction-seconds of static injection, and all four paced-minus-static intervals include zero. Bad requests per fraction-second stayed at about 30--34 at 32 requests/s and 62--66 at 64 requests/s, matching Table~\ref{tbl:cost}, so the extra exposure cost the expected number of bad requests.

The gain came with a small loss of safety margin. Nine of the 20 runs exceeded the usable guard by 1--4 requests, and seven of those were ended early by the common allocation safety guard; the other 13 ran the full window. Mean guard overshoot was 0.8--1.6 requests, compared with 4.6--21.6 for static injection (paired reductions of 3.6--20.0 requests) and almost zero for the closed-loop policy. The 675- and 1,350-request reserves absorbed every overshoot: the largest total was 964 of a 2,310-request allocation. Pacing against the remaining allocation therefore removes the burn-gate under-use. However, it moves the policy close to the guard, and its residual overshoot comes from guard trips once the remaining allocation is nearly spent. A randomized study that runs paced and closed-loop control in the same blocks is still needed before claiming the gain causally.

\begin{table*}[width=\textwidth,cols=7,pos=t]
\caption{Budget-paced follow-up over five blocks. Guard overshoot is in requests and exposure in fraction-seconds, on the latency ledger (an abort counts against both ledgers, so abort values are the same for availability). Means for paced runs, and paced minus v7 differences with 95\% block-bootstrap intervals. The pairing is time-separated, not jointly randomized.}\label{tbl:followup}
\footnotesize
\begin{tabular*}{\tblwidth}{@{\extracolsep{\fill}}llrrrll@{}}
\toprule
Fault & Req/s & Guard overshoot & Exposure & Usable spent & $\Delta$ exposure vs.\ closed-loop & $\Delta$ exposure vs.\ open-loop \\
\midrule
abort & 32 & 1.0 & 15.41 & 98\% & $10.95$ [$10.47$, $11.39$] & $2.00$ [$1.35$, $2.64$] \\
abort & 64 & 1.6 & 15.09 & 100\% & $10.88$ [$10.60$, $11.05$] & $1.69$ [$1.35$, $2.03$] \\
delay & 32 & 1.0 & 15.03 & 99\% & $0.98$ [$0.06$, $1.59$] & $1.80$ [$1.15$, $2.35$] \\
delay & 64 & 0.8 & 14.81 & 99\% & $0.76$ [$0.13$, $1.49$] & $1.56$ [$1.30$, $1.87$] \\
\bottomrule
\end{tabular*}
\end{table*}

